\section{Machine Learning in Product Development and Production}
\label{sec:ml-in-product-development-and-production}

This section reviews how machine learning is applied in industrial
product development and production. It first describes common
application areas, then the learning paradigms used in these
settings. The last part summarises the findings of the pre-master
project that this thesis builds on.

\subsection{Applications of Machine Learning}
\label{sec:applications-of-machine-learning}

Machine learning is widely applied in industrial product development and production, 
where it supports tasks such as process optimisation, quality prediction, fault diagnosis, 
predictive maintenance, and anomaly detection 
\parencite{wuest2016ml_manufacturing, plathottam2023review_ai_manufacturing}. 
The application of ML in these domains continues to grow, 
driven by the increasing availability of production data and the improved usability of ML tools 
\parencite{wuest2016ml_manufacturing}. 
Unlike traditional rule-based approaches, ML methods can learn directly from data 
and adapt to changing conditions, making them well suited to the complex and dynamic 
nature of manufacturing environments \parencite{wuest2016ml_manufacturing}. 
In practice, ML is used both to improve product quality and to optimise production processes, 
for example through early prediction of manufacturing outcomes, root cause analysis, 
and condition monitoring of equipment \parencite{weichert2019review_ml_optimization_production}. 
As well as individual processes, ML also supports broader operational tasks such as 
production planning, scheduling, and supply chain management 
\parencite{plathottam2023review_ai_manufacturing, presciuttini2024ml_iot_manufacturing_operations}.

\subsection{Learning Paradigms in Industrial Applications}
\label{sec:learning-paradigms-in-industrial-applications}

Across reported applications in manufacturing and production, supervised learning 
is the dominant paradigm, accounting for the large majority of studies in most 
application areas \parencite{wuest2016ml_manufacturing, adekoya2022applications_ai_em}. 
This reflects the nature of many industrial problems, where labelled data are available 
through quality inspections, sensor logs, and operator feedback, and where the goal 
is typically to predict a known outcome such as product quality or equipment failure 
\parencite{wuest2016ml_manufacturing, presciuttini2024ml_iot_manufacturing_operations}. 
Unsupervised learning plays a smaller but important role, particularly in anomaly 
detection and condition monitoring, where labelled examples of faults are scarce or 
costly to obtain \parencite{presciuttini2024ml_iot_manufacturing_operations}. 
Reinforcement learning is less commonly reported in industrial 
settings \parencite{wuest2016ml_manufacturing}, though it has shown 
strong results in process control and scheduling problems, where 
decisions must be made sequentially over time 
\parencite{panzer2022deep_reinforcement_learning_production_systems}.
In practice, the choice of learning paradigm is often shaped by what data is available 
rather than by the structure of the problem itself 
\parencite{wuest2016ml_manufacturing}.

\subsection{Pre-Master Findings}
\label{sec:pre-master-findings}

In a preceding specialisation project \parencite{bergsto2026mlselection}, referred to as the pre-master 
project, a large-scale analysis of machine learning 
research in product development and production was conducted, classifying over 33,000 
article abstracts across product lifecycle phase, application area, learning paradigm, 
and reported ML methods. 
The results showed that the literature is strongly concentrated in the optimisation
phase of the product lifecycle, while early phases such as planning and concept 
development contain comparatively few studies. 
The analysis confirmed the paradigm distribution described in 
Section~\ref{sec:learning-paradigms-in-industrial-applications}: 
supervised learning dominated across all phases and application areas, 
accounting for approximately 73--80\% of the articles in each lifecycle 
phase, with neural networks, random forests, and support vector machines 
as the most frequently reported methods. Unsupervised learning appeared 
mainly in anomaly detection and predictive maintenance, while 
reinforcement learning was the least frequent paradigm and was largely 
limited to control and scheduling problems.